\documentclass[11pt,a4paper]{article}
\usepackage[margin=1in]{geometry}
\usepackage{amsmath,amssymb,booktabs}
\usepackage[hidelinks]{hyperref}
\raggedbottom
\widowpenalty=10000
\clubpenalty=10000
\setlength{\emergencystretch}{2em}
\newcommand{\Mpl}{M_{\rm Pl}}
\newcommand{\GN}{G_N}
\begin{document}
\title{A Terrestrial Test of Radial Dimensionality Theory: Source-Mass Density Dependence of Newton's Constant}
\author{Christopher K. Merrill\\\small Independent Researcher\\[1ex]
Claude (Anthropic)\\\small Large Language Model, Computational Collaboration\\[1ex]
Astra (Open AI)\\\small Large Language Model, Computational Collaboration}
\date{Consistency revision: November 12, 2026}
\maketitle
\begin{abstract}
We separate a proposed density-dependent Poisson law from Paper VI's
scalar action, whose claimed equilibrium fails a source-sign check. The
Poisson ansatz still defines conditional laboratory force predictions.
We correct the uncertainty units in a retained NIST comparison and propagate
the published four-result covariance. We withdraw the cross-experiment
amplitude thresholds: the assigned single density for a mixed-material
consensus, apparatus confounding and interval construction do not support
those bounds. The resulting material contrasts are diagnostics, not an RDT
detection or a calibrated amplitude exclusion.
\end{abstract}

\section*{Version note: v3}
This version retains the published title and a testable laboratory proposal. Source-density dependence is a separately specified phenomenological force-law ansatz, not a derivation from the failed action in Paper VI. The revision corrects NIST uncertainty units and covariance propagation and withdraws unsupported cross-experiment amplitude bounds. Retained measurement contrasts are descriptive diagnostics, not a detection or calibrated bound.


\section{Specified phenomenological force model}
Adopt, as a separate hypothesis,
\begin{equation}
 \nabla^2\Phi=4\pi\GN\frac{\rho}{1-\sigma(\rho)},\qquad
 \sigma(\rho)=A\frac{x}{1+x},\quad x=(\rho/\rho_0)^\alpha.
\end{equation}
For an isolated uniform spherical source of mass $M$ and radius $R$, with
ordinary test acceleration $-\nabla\Phi$, integration gives
\begin{equation}
 \Phi'(r)=\frac{\GN M}{[1-\sigma(\rho_s)]r^2},\qquad r>R.
\end{equation}
Spherical symmetry is required for this inverse-square monopole form.
An arbitrary apparatus needs its geometry integrated; an inhomogeneous body
has weighted source charge $\int \rho/[1-\sigma(\rho)]\,d^3x$.
Zero local exterior density does not erase this integrated charge.
The exact two-uniform-material ratio and its first-order approximation are
\begin{equation}
 D_{12}=\frac{G_1}{G_2}-1
 =\frac{\sigma_1-\sigma_2}{1-\sigma_1}
 =\sigma_1-\sigma_2+O(A^2).
\end{equation}
Only the leading expression scales exactly linearly with $A$.
For illustrative densities 19.3 and $2.7\,{\rm g\,cm^{-3}}$, $\rho_0=150$
and fixed $A=0.0087$, the retained script gives:
\begin{center}
\begin{tabular}{lll}
\toprule
$\alpha$ & Leading (ppm) & Exact (ppm) \\
\midrule
0.02 & 85.48 & 85.84 \\
0.12 & 496.19 & 498.09 \\
0.334 & 1113.26 & 1116.51 \\
0.5 & 1267.67 & 1270.59 \\
\bottomrule
\end{tabular}

\end{center}
These are chosen model inputs, not measured material properties or points
on the alpha-dependent solar threshold.

The ansatz specifies an active-source force calculation. It does not supply
a self-consistent passive/inertial mass law, scalar surface matching or a
momentum-conserving two-body action. For unequal density factors, blindly
assigning ordinary passive masses to both bodies gives unequal mutual forces;
additional physics would be needed to close that model. A source/test-body
configuration can still test the stated force prescription, with its scope
explicit. It is not derived by Paper VI's audited action or by the full RDT
operator.

\section{Solar comparison and model distinction}
The canonical operator-model scan~\cite{Series_paper1} uses $\rho_0=150\,{\rm g\,cm^{-3}}$
and $\alpha\in[0.02,0.5]$. Its nominal solar proxy threshold, denoted
$A_{\rm lim}$, has a sampled and locally refined amplitude peak
$A=0.0087502019$ at $\alpha=0.1244971139$ (asymptote $3-3A=2.973749394$).
It has neither calibrated confidence coverage nor a certified continuous
maximum. Changing the physical response requires a new solar comparison;
sharing a density-law shape does not transfer this threshold.

In particular, an illustrative $A=0.0087$ is not a bound on the terrestrial
ansatz. Paper VI's specified increasing potential~\cite{Series_paper6} fails its positive-density
adiabatic equilibrium condition, so its former screening argument cannot
justify neglecting apparatus or surface effects here. The old scalar-tensor
companion~\cite{Series_companion} is another distinct construction. This note does not repair either
action by fitting laboratory data.

\section{Published NIST table and corrected uncertainties}
The final published PDF of Ref.~\cite{NIST}, Table 16, gives the following
four results. The uncertainty column is relative standard uncertainty in
ppm; the former notation such as $6.673642(23)$ incorrectly treated ppm
as last-digit errors.
\begin{center}
\begin{tabular}{lrr}
\toprule
Configuration & $G/(10^{-11}\,{\rm m^3\,kg^{-1}\,s^{-2}})$ & $u_r$ (ppm)\\
\midrule
Copper, servo & 6.673642 & 23\\
Copper, free & 6.674021 & 30\\
Sapphire, servo & 6.672637 & 38\\
Sapphire, free & 6.673636 & 94\\
\bottomrule
\end{tabular}
\end{center}
Table 18 provides the less-rounded uncertainties
$(23.2,30.3,37.5,93.9)$ ppm and correlation matrix
\begin{equation}
 C=\begin{pmatrix}
 1&.42&.38&.12\\ .42&1&.23&.23\\ .38&.23&1&.25\\ .12&.23&.25&1
 \end{pmatrix}.
\end{equation}
The published consensus is $6.67387\pm0.00038$ in the same units.
It combines materials and has no unique source density. We use the final
PDF rather than the conflicting NIST landing-page abstract.

For contrast vector $c$, set $\Sigma_{ij}=u_i u_jC_{ij}$ and calculate
$c^TG$ and $\sqrt{c^T\Sigma c}$. Dividing both by a fixed reference
$G_{\rm ref}=6.6737\times10^{-11}$ gives the following descriptive ppm
contrasts; this normalization is not a fit parameter.
\begin{center}
\begin{tabular}{lll}
\toprule
Contrast & Difference (ppm) & $u$ (ppm) \\
\midrule
Cu--Sa, servo & 150.59 & 35.82 \\
Cu--Sa, free & 57.69 & 91.80 \\
Free--servo, Cu & 56.79 & 29.43 \\
Free--servo, Sa & 149.69 & 91.99 \\
\bottomrule
\end{tabular}

\end{center}
These uncertainties propagate the published measurement budget only; they
do not include fitted configuration biases or the authors' additional dark
uncertainty. They must not be read as detection significances. Both material
contrasts are positive, but fixed-material method contrasts are comparable.
The source discusses unresolved offsets and configuration-specific biases;
it does not identify their cause conclusively. The earlier claim that it
established composition was not the cause is withdrawn.

\section{Why the earlier regression is not an amplitude bound}
The preserved \texttt{g\_regression\_v2.py} fits
$G_i=a+b\,x_i/(1+x_i)$, reports $A=b/a$, inflates an approximate error by
a constant-model Birge ratio and forms
$\max(\widehat A,0)+1.645u_A$. This is an exploratory recipe, not a
coverage-calibrated construction at a physical boundary. Its $u_A$ also
omits intercept uncertainty and slope--intercept covariance in the ratio.
A delta-method variance would contain
\begin{equation}
 \operatorname{Var}(b/a)\simeq
 \frac{\operatorname{Var}b}{a^2}
 +\frac{b^2\operatorname{Var}a}{a^4}
 -\frac{2b\operatorname{Cov}(a,b)}{a^3}.
\end{equation}
Fixing that term alone would not solve the apparatus correlations or
material-dependent biases. The source-density assignment for the mixed
NIST consensus is invalid even if one configuration carries more statistical
weight. A new regression would need a response model for the actual mixture.

The historical 12-row dataset and regression scripts are retained for
provenance; the other experiments' source materials and revisions have not
been re-audited in this step. Their fitted values and the ``stronger than
solar'' comparison are therefore removed from the current evidence tables.
No replacement multi-experiment inference is claimed. The retained four-row
source check is identified explicitly rather than called a verification of
all historical data.

\section{Empirical next steps and conclusion}
A useful force test specifies source and test-body geometry, measured material
properties, calibration, environmental fields and the covariance of the
measurements. A material swap can reduce apparatus differences but does not
by itself isolate density: source mass, surface properties, temperature and
method-dependent torques can also change. A model for these effects can be
compared directly with data; outside endorsement is not a prerequisite.
The present audit corrects arithmetic and provenance and leaves the
phenomenological force hypothesis testable. It establishes neither a detection,
a calibrated exclusion nor a scalar completion of the operator theory.
\section*{Code and revision availability}
Producing scripts, inputs, saved results and numerical audit evidence are
publicly available in the \href{https://github.com/chrismerrill1974/solarcore_code/tree/main/reproduction}{shared RDT reproducibility archive}.
The accompanying source package, \nolinkurl{merrill2026_terrestrial_g_code}, includes the
manuscript, dependencies and resource-controlled reproduction commands.
Its \texttt{reproduction/CAPSULE.json} records file hashes;
\texttt{reproduce.py} verifies the package and rebuilds the paper from retained
scientific outputs. Full scientific runs are separate from this saved-result
reproduction. Historical scripts and results remain separately identified.

\section*{Acknowledgments}
Developed with computational collaboration from Claude (Anthropic) and
Astra (Open AI). Earlier claims and calculations remain identifiable in the
repository history; this revision supersedes their interpretation.
\begin{thebibliography}{9}
\bibitem{NIST} S.~Schlamminger et al., ``Redetermination of the gravitational
constant with the BIPM torsion balance at NIST,'' Metrologia
\textbf{63}, 025012 (2026),
\href{https://doi.org/10.1088/1681-7575/ae570f}{doi:10.1088/1681-7575/ae570f},
Tables 16 and 18.
\bibitem{Series_companion} C.~K.~Merrill, Claude (Anthropic), and Astra (Open AI),
``A Scalar-Tensor Companion to Radial Dimensionality Theory: A Chameleon Model Reproducing the Density Law, with Gravitational-Wave Predictions,''
revised version, November 12, 2026.

\bibitem{Series_paper1} C.~K.~Merrill, Claude (Anthropic), and Astra (Open AI),
``Radial Dimensionality Theory I: Solar Neutrino and Helioseismic Constraints on Density-Dependent Spatial Dimensionality,''
revised version, November 12, 2026.

\bibitem{Series_paper6} C.~K.~Merrill, Claude (Anthropic), and Astra (Open AI),
``Radial Dimensionality Theory VI: A Field-Theoretic Origin for the Density-Dependent Gravitational Coupling,''
revised version, November 12, 2026.

\end{thebibliography}
\end{document}
